\section*{3.2 Weakly resistant gadgets (author source pp.16--17)}
There are gadgets, such as $x\vee(y\wedge z)$, which are clearly not resistant, but for which the lifting still holds as we will show below. To capture these cases, we define a more general class of weakly resistant gadgets.
\begin{originalresult}{Definition 3.3}
A gadget $g:\bits^m\to\bits$ is weakly resistant if for every certificate $\alpha$ (a partial assignment which sets the value of $g$) there is a substitution of one input variable (with 0 or 1) which conflicts with $\alpha$ and does not make $g$ constant.
\end{originalresult}
Every resistant gadget is trivially weakly resistant: we can take any variable mentioned in $\alpha$ and substitute the opposite value. The aforementioned gadget $h(x,y,z)=x\vee(y\wedge z)$ is weakly resistant (as we show below) but not resistant, since $h|_{x=1}\equiv1$. The gadget $x\vee y$ is not even weakly resistant due to the certificate $x=y=0$: if we substitute $x=1$ or $y=1$ then the gadget becomes constant.
Let us verify that $h$ is weakly resistant by considering all possible certificates:
\begin{itemize}
\item $x=1$: take $x=0$.
\item $y=z=1$: take $y=0$ or $z=0$.
\item $x=y=0$: take $y=1$.
\item $x=z=0$: take $z=1$.
\end{itemize}
We now prove the following lifting theorem:
\begin{originalresult}{Theorem 3.4}
For any $f:\bits^n\to\bits$ and a weakly resistant gadget $g:\bits^m\to\bits$ the following holds:
\[\log\DTSize(f\comp g)=\Omega(\CC(f)).\]
\end{originalresult}
\begin{proof}
Given a decision tree $T$ for $f\comp g$ of size $S$, we construct a blockwise random projection $\rho$ in the following way. For each $i\in[n]$, we choose a random certificate $C_i$ of $g$. Let $z_{i,j}=b$ be the substitution promised by weak resistance for $C_i$ that does not make $g(z_i)$ constant. We can complete $\rho|_{z_i}$ to a projection such that $g(\rho|_{z_i})=x_i$ (as in the case of resistant predicates, see proof of Theorem 3.2 using random projections).
Let $\ell$ be a leaf of $T$ and $\tau$ be the corresponding partial assignment to the variables of $f\comp g$. Suppose that $\tau$ contains certificates for at least $D$ blocks (i.e., in these blocks, the value of $g$ on $\tau$ is fixed). If there are $M$ certificates of $g$, then $\tau$ is consistent with $\rho$ with probability at most $(1-1/M)^D$. Choose $D=\Theta(\log S)$ so that $(1-1/M)^DS<1$. By union bound, under some choice of $\rho$, the partial assignments for all leaves of $T|_\rho$ contain certificates for less than $D$ blocks. Since by construction $(f\comp g)|_\rho=f$, this shows that $\CC(f)<D$.
Indeed, given an input $x$ to $f$, consider the corresponding input $\rho(x)$ to $f\comp g$. This input reaches some leaf $\ell$, in which less than $D$ blocks contain certificates, say the blocks in $I\subseteq[n]$. For each input $x'$ consistent with gadget values for blocks in $I$, we can find an input $z$ reaching the same leaf such that $g^n(z)=x'$. Hence, the block certificates in $\ell$ correspond to a certificate of $f$.
\end{proof}
\paragraph{Remark.} One of the anonymous reviewers noticed that the structure of a decision tree is completely irrelevant to the above proof. We only use the fact that the leaves are monochromatic and cover all the inputs. Thus, Theorem 3.4 can be restated by replacing the decision tree size with the subcube cover number. This also explains why, in the next section, we will be able to reuse the above proof in a similar lifting result for conjunction DAGs (see Theorem 4.2).
\begin{originalresult}{Corollary 3.5}
For any $f:\bits^n\to\bits$ and a weakly resistant gadget $g:\bits^m\to\bits$ the following holds:
\[\log\DTSize(f\comp g)=\Omega(\sqrt{\DT(f)}).\]
\end{originalresult}
\begin{proof}
Beals et al. (2001) shows that for any function $f$,
\[\DT(f)\le(\CC(f))^2.\]
Together with Theorem 3.4, that gives us the desired bound.
\end{proof}
